\begin{theorem}[The two minimal smooth fillings]\hypertarget{tag-02ZK}{}\label{push:src:A.2.2}\label{push:A22}
\StatementRecord{push:src:A.2.2}{proved}{02ZK}{\SeriesRef{push:src:A.2.1}, \SeriesRef{push:A12}, \SeriesRef{push:src:A.1.2}}

The two minimal subdivisions of \ref{push:src:A.2.1} give proper flat holomorphic
families $\mathcal X^+,\mathcal X^-\to\Delta_u$ with smooth total
spaces. Their joint maps to $B\times\Delta_u$ are flat. Over
$\Delta_u^*$ they are isomorphic as families to the specified
untwisted family $X(u)$. Each central fibre $Q^\pm$ is reduced,
irreducible, and has normal crossings, with smooth double locus
$S$ and no triple points.

Its normalization $R^\pm$ is obtained by blowing up $R$ along
$S_\infty$ in, respectively, one of its two sections and then
taking the small resolution at the unique ordinary double point
of that blow-up prescribed by the fixed bottom diagonal.
The two proper transforms of the sections are still copies of
$S$, glued by $T$. Their normal-bundle product is trivial;
in particular $\mathcal T^1_{Q^\pm}\simeq\mathcal O_S$.
\end{theorem}


\begin{proof}
\begin{enumerate}[label=\arabic*.,leftmargin=*]
\item Apply either fan subdivision over the cusp neighbourhood in
Theorem~\ref{push:A12}, and use the identity modification elsewhere.
Toric modifications are proper, and a subdivision is an
isomorphism off the strata it subdivides. Its restriction over
$q\ne0$ is the identity, so this gluing has no additional
transition choice. The period actions preserve the subdivisions,
and the covering description in \ref{push:src:A.1.2} descends the proper
modifications to the quotient. All maximal cones are regular.
Every local ring of the resulting fourfold is therefore regular.
The coordinate $u$ is the product of those affine coordinates
whose rays have $h_u=1$; there are at most two such rays in a
maximal cone. This proves reduced normal crossings and the absence
of triple points for the $u$-central fibre. The regular-sequence
argument, or the saturated face images in \ref{push:src:A.2.1}, proves joint
flatness. Properness follows from properness of the toric
modification and of $\mathcal X^\flat$.

\item We give the normalization calculation, including the singular
centre. At a node of $S_\infty$, write the original normalization
as coordinates $(X,z,v)$, with section $X=0$ and centre
$(X,zv)$. Blowing up that ideal has charts
\[
 X=zvA;\qquad zv=XB.
\]
The first is smooth, and the second has one conifold at its
origin. A small resolution of the latter is, for one of the
bottom-diagonal conventions,
\[
 X=za,\quad v=aB;\qquad z=Xd,\quad B=dv.
\]
Interchanging $z,v$ gives the other diagonal convention. These
are holomorphic charts of a proper small modification: they are
the two usual charts of the blow-up of the appropriate rank-one
ideal on the conifold. The strict transform of $X=0$ is $A=0$
in the first chart and never passes through the conifold. It is
therefore unchanged as a surface. The other section is disjoint
from the centre and is unchanged as well. Over a smooth point
of $S_\infty$ the exceptional fibre is one $\mathbb P^1$;
over its node it is a chain of two $\mathbb P^1$'s, one being
the exceptional curve of the small resolution. This identifies
the normalization morphism, including its order of operations.

\item For example the corresponding total fourfold modification of
$xy=uzv$ can be written by first taking $x=ua$, $ay=zv$ and
then resolving the conifold in $(a,y,z,v)$, together with the
other chart $u=xb$, $y=bzv$. Its three smooth affine charts map
to $(x,y,z,v,u;q)$ by
\[
\begin{aligned}
 (u,z,\lambda,y)&\longmapsto(uz\lambda,y,z,\lambda y,u;z\lambda y),\\
 (u,a,\nu,v)&\longmapsto(ua,\nu v,a\nu,v,u;a\nu v),\\
 (x,b,z,v)&\longmapsto(x,bzv,z,v,xb;zv).
\end{aligned}
\]
The first two are the conifold-resolution charts
$a=z\lambda$, $v=\lambda y$ and $z=a\nu$, $y=\nu v$.
Substitution gives $xy=uzv$ and $q=zv$ in every chart.
Exchanging the chosen zero and infinity branches gives the other model.

\item The normal bundle of the section containing the centre changes
by tensoring with $\mathcal O_S(-S_\infty)$; the second normal
bundle is unchanged. Since the original product is $J$, the
new product is $J(-S_\infty)=\mathcal O_S$. The local
formula for $\mathcal T^1$ is the product of the two normal
bundles, with $T$ used on the infinity branch. This proves
d-semistability in this precise sense. The $\rho_l$ form one
orbit, so the central fibre is irreducible and its conductor is
exactly the two stated sections. Over $u\ne0$ both subdivisions
restrict to the fixed bottom triangulation, giving precisely
Engel's Mumford filling, simultaneously in $u$.

\item The birational comparison $\mathcal X^+\dashrightarrow
\mathcal X^-$ is the identity over the punctured $u$ disc and
off the exceptional loci over $(q,u)=(0,0)$. Locally it replaces
the two compatible side-face diagonals by the other two.
All its rays lie in the Gorenstein hyperplane $h_q+h_u=1$,
so both small modifications are crepant over
$\mathcal X^\flat$. Over a smooth point of the nodal curve
this is the ordinary Atiyah flop times a curve; at the node its
extension is exactly the displayed pair of prism subdivisions.
This specifies the birational map.
\end{enumerate}
\end{proof}
